\CourseSection{Risks and Safety Hazards}\label{sec:risks}
% TODO (handoffs from Section 4):
% - Speed chain (resolved): t_r = 0.3 s (proposal PF-01) and a = 2 m/s^2 (MC-06)
%   give a 1.0 m stop at 1.5 m/s, so MC-08 = 1.5 m/s, FL-03 min = 1 m,
%   OB-03 = 1 m, FL-04 = 0.5 s (0.75 m blind travel). MC-03 = 1.4 m/s (Bohannon).
% - Team: agree on the proposed initial/residual ratings.
% - TODO: verify the matrix cell mapping against MIL-STD-882E Table III.
% - TODO: the 0.7 m arm length is an uncited estimate. gordon2014ansur (ANSUR II) is in
%   references.bib if a slot opens, but confirm the value first.

\CourseSubsection{Hazard Identification and Risk Assessment}\label{sec:risks-assessment}

Hazards were identified and reduced using the process in ISO~12100~\cite{iso12100} and rated with the MIL-STD-882E risk matrix~\cite{milstd882e}, which scores severity from 1 (catastrophic) to 4 (negligible) and likelihood from A (frequent) to E (improbable). Table~\ref{tab:hazards} (Appendix~\ref{app:hazards}) lists each hazard, its cause, the standard that covers it, its size, and its rating before and after mitigation. The most serious is the cart moving without a valid command, for example after the Pi crashes. Most mitigations make a hazard less likely rather than less harmful, and the matrix rates a moderate-harm hazard as medium whether it is occasional or improbable, so most ratings stay medium. Technical and schedule risks are covered in Section~\ref{sec:approach-contingency}.

Speed sets the size of the motion hazards. A cart moving at speed $v$ that reacts after time $t_r$ and brakes at $a$ needs a stopping distance of
\begin{equation}
d = v\,t_r + \frac{v^2}{2a} .
\label{eq:stopping}
\end{equation}
With a 0.3\,s reaction time and 2\,m/s$^2$ braking, the cart stops within 1.0\,m at 1.5\,m/s but needs 1.6\,m at 2\,m/s. This is why speed is capped at 1.5\,m/s and the cart keeps at least 1\,m from the golfer. The drive test (DRV-4) measures the actual braking on grass, and the speed cap is lowered if it falls short (C-9).

\CourseSubsection{Risk Mitigation and Residual Risk}\label{sec:risks-mitigation}

Table~\ref{tab:mitigations} and Figure~\ref{fig:bowtie} (Appendix~\ref{app:hazards}) give each mitigation, the requirement it creates, and how it is verified. The key choice is where the stops are enforced. The speed and current limits, and the stop when commands from the Pi stop arriving, run in the mainboard firmware, so they hold even if the Pi crashes. The e-stop first commands the motors to brake, then a relay cuts motor power without going through the Pi, because cutting power alone would let the motors coast. This is a Category~1 stop~\cite{iec60204-1}, and the cart does not restart until it is manually reset~\cite{iso13850}.

The remaining risks are accepted because a safety operator holds the e-stop during every test in which the cart moves, on flat ground in dry weather. Four become limitations in Section~\ref{sec:conclusion}: the cart can roll on a slope once power is cut, the sensors can miss very thin or low objects, a battery fire would still be severe, and the golfer may swing near a parked cart.
